\subsection{Additional Quantitative Analysis of Main Results} \label{appx:sub:quantitative}
We supplement Table~\ref{tab:main} with fine-grained rollout success of a few baseline methods compared to the \RLM(recursion depth=1). In Figure~\ref{fig:tab1_comparison}, we generally find \RLM{}s solve the same tasks, and more tasks than, other baselines, especially for GPT-5.

\begin{figure}[!hbt]
    \centering
    \includegraphics[width=1\linewidth]{figures/table_1_comparison.png}
    \caption{For GPT-5 and Qwen3-Coder-480B-A35B-Instruct, we plot how the \RLM(recursion depth=1) compares to other baselines using the same models. For all tasks where at least one method gets the answer correct, we show how many only the \RLM{} got correct in green, how many both got correct in gray, and how many only the baseline got correct in red.}
    \label{fig:tab1_comparison}
\end{figure}

We also explore how sub-calling behavior differs between rollouts. In Figure~\ref{fig:sub-call-rollouts}, we find a wide range of sub-calling behaviors that greatly differ across models and even for correct and incorrect rollouts. For example, GPT-5 uses significantly more sub-calls for BrowseComp-Plus than any other model. However, for OOLONG, Qwen3-Coder uses a large number of sub-calls (~500 on average) for correct rollouts, which is significantly more than the number used by GPT-5. Furthermore, Qwen3-8B in particular tends to use more sub-calls on incorrect trajectories.

\begin{figure}[!hbt]
    \centering
    \includegraphics[width=1\linewidth]{figures/sub-call-rollouts.png}
    \caption{For each task in Table~\ref{tab:main}, we plot the average number of sub-calls made during an \RLM{}(depth=1) trajectory for each task, grouped by whether it got the task correct or incorrect.}
    \label{fig:sub-call-rollouts}
\end{figure}

\subsection{Additional Runtime and Cost Analysis of \RLM{}s}
We supplement the cost and runtime analysis of \RLM{}s with additional, fine-grained plots. We focus on RLMs with depth=0 (i.e. no sub-calls) and depth=1. In Figures~\ref{fig:cost-gpt-5}, \ref{fig:cost-qwen3} we include a histogram for the cost of each method on every task for both GPT-5 and Qwen3-Coder. We generally observe long-tailed, high-variance trajectories for \RLM{}s in both models. We plot the cost of \RLM(depth=1) and baselines at quartiles in Figure~\ref{fig:quartiles}.

We additionally include log-scaled runtime plots (Figure~\ref{fig:runtime-gpt-5},~\ref{fig:runtime-qwen3}) for each method below. The tail end (e.g. 95th percentile) shows extremely long runtimes, which is mainly due to sequential sub-LLM calls taking up most of the runtime. However, we observe these cases happen infrequently, and can be early-stopped with timeout logic. As we remarked in \S\ref{sec4.4-qualitative}, the runtime for these methods can be significantly improved through asynchrony of LM calls and additional prompting to discourage long sub-LM calls or code.

For the scaling plot in Figure~\ref{fig:rlm-scaling}, we also provide the average API cost per task.

\begin{figure*}[htb!]
    \centering
    \includegraphics[width=\textwidth]{figures/cost_quartiles_dual_depth_with_baselines.png}
    \caption{Cost of \RLM{} and baselines described in \S\ref{sec4.2-methods} plotted at the 25th, 50th, 75th, and 95th percentile of total API cost. We observe comparable or even lower costs for \RLM{}s at the 50th percentile, but sharp increases at the tail end due to potentially long \RLM{} trajectories.}
    \label{fig:quartiles}
\end{figure*}

\begin{figure}[htb!]
    \centering
    \includegraphics[width=0.9\textwidth]{figures/runtime_quartiles_gpt-5.png}
    \caption{Plotted quartiles of the runtime for methods and baselines around GPT-5 across OOLONG, OOLONG-Pairs, CodeQA, and BrowseComp+ (1K) for all methods described in \S\ref{sec4.2-methods}. We plot the 25th, 50th, 75th, and 95th percentiles.}
    \label{fig:runtime-gpt-5}
\end{figure}

\begin{figure}%
    \centering
    \includegraphics[width=0.9\textwidth]{figures/runtime_quartiles_qwen3-coder-480b-a35b-instruct.png}
    \caption{Plotted quartiles of the runtime for methods and baselines around Qwen3-Coder-480B-A35B-Instruct across OOLONG, OOLONG-Pairs, CodeQA, and BrowseComp+ (1K) for all methods described in \S\ref{sec4.2-methods}. We plot the 25th, 50th, 75th, and 95th percentiles.}
    \label{fig:runtime-qwen3}
\end{figure}

\begin{figure}%
    \centering
    \includegraphics[width=1\textwidth]{figures/cost_distributions_gpt-5.png}
    \caption{Histogram of the API costs for GPT-5 across OOLONG, OOLONG-Pairs, CodeQA, and BrowseComp+ (1K) for all methods described in \S\ref{sec4.2-methods}.}
    \label{fig:cost-gpt-5}
\end{figure}

\begin{figure}%
    \centering
    \includegraphics[width=1\textwidth]{figures/cost_distributions_qwen3-coder-480b-a35b-instruct.png}
    \caption{Histogram of the API costs for Qwen3-Coder-480B across OOLONG, OOLONG-Pairs, CodeQA, and BrowseComp+ (1K) for all methods described in \S\ref{sec4.2-methods}.}
    \label{fig:cost-qwen3}
\end{figure}

\begin{figure}%
    \centering
    \includegraphics[width=\textwidth]{figures/scaling_cost.png}
    \caption{We plot the API cost in USD for the runs in Figure~\ref{fig:rlm-scaling}.}
    \label{fig:cost-scaling}
\end{figure}
